\documentclass[10pt,a4paper]{amsart}
\usepackage{amsmath,amssymb,amsthm,graphicx}
\usepackage{verbatim}
\newenvironment{reference}{\comment}{\endcomment}
%\newenvironment{reference}{}{}
\newenvironment{slogan}{\comment}{\endcomment}
\newenvironment{history}{\comment}{\endcomment}

% For commutative diagrams we use Xy-pic
\usepackage[all]{xy}

% We use 2cell for 2-commutative diagrams.
\xyoption{2cell}
% Theorem environments.
%
\theoremstyle{plain}
\newtheorem{theorem}[subsection]{Theorem}
\newtheorem{proposition}[subsection]{Proposition}
\newtheorem{lemma}[subsection]{Lemma}

\theoremstyle{definition}
\newtheorem{definition}[subsection]{Definition}
\newtheorem{example}[subsection]{Example}
\newtheorem{exercise}[subsection]{Exercise}
\newtheorem{situation}[subsection]{Situation}

\theoremstyle{remark}
\newtheorem{remark}[subsection]{Remark}
\newtheorem{remarks}[subsection]{Remarks}

\numberwithin{equation}{subsection}

% Macros
%
\def\lim{\mathop{\mathrm{lim}}\nolimits}
\def\colim{\mathop{\mathrm{colim}}\nolimits}
\def\Spec{\mathop{\mathrm{Spec}}}
\def\Hom{\mathop{\mathrm{Hom}}\nolimits}
\def\Ext{\mathop{\mathrm{Ext}}\nolimits}
\def\SheafHom{\mathop{\mathcal{H}\!\mathit{om}}\nolimits}
\def\SheafExt{\mathop{\mathcal{E}\!\mathit{xt}}\nolimits}
\def\Sch{\mathit{Sch}}
\def\Mor{\mathop{\mathrm{Mor}}\nolimits}
\def\Ob{\mathop{\mathrm{Ob}}\nolimits}
\def\Sh{\mathop{\mathit{Sh}}\nolimits}
\def\NL{\mathop{N\!L}\nolimits}
\def\CH{\mathop{\mathrm{CH}}\nolimits}
\def\proetale{{pro\text{-}\acute{e}tale}}
\def\etale{{\acute{e}tale}}
\def\QCoh{\mathit{QCoh}}
\def\Ker{\mathop{\mathrm{Ker}}}
\def\Im{\mathop{\mathrm{Im}}}
\def\Coker{\mathop{\mathrm{Coker}}}
\def\Coim{\mathop{\mathrm{Coim}}}

% Boxtimes
%
\DeclareMathSymbol{\boxtimes}{\mathbin}{AMSa}{"02}

%
% Macros for moduli stacks/spaces
%
\def\QCohstack{\mathcal{QC}\!\mathit{oh}}
\def\Cohstack{\mathcal{C}\!\mathit{oh}}
\def\Spacesstack{\mathcal{S}\!\mathit{paces}}
\def\Quotfunctor{\mathrm{Quot}}
\def\Hilbfunctor{\mathrm{Hilb}}
\def\Curvesstack{\mathcal{C}\!\mathit{urves}}
\def\Polarizedstack{\mathcal{P}\!\mathit{olarized}}
\def\Complexesstack{\mathcal{C}\!\mathit{omplexes}}
% \Pic is the operator that assigns to X its picard group, usage \Pic(X)
% \Picardstack_{X/B} denotes the Picard stack of X over B
% \Picardfunctor_{X/B} denotes the Picard functor of X over B
\def\Pic{\mathop{\mathrm{Pic}}\nolimits}
\def\Picardstack{\mathcal{P}\!\mathit{ic}}
\def\Picardfunctor{\mathrm{Pic}}
\def\Deformationcategory{\mathcal{D}\!\mathit{ef}}
\usepackage{cleveref}
\crefname{originalreference}{original reference}{original references}
\Crefname{originalreference}{Original reference}{Original references}
\makeatletter
\newlabel{lemma-node-over-separable-extension}{{0CD6}{source}{}{sourcepointer-1}{}}
\newlabel{lemma-node-over-separable-extension@cref}{{[originalreference][0][]0CD6}{[1][0][]source}{}{}{}}
\makeatother
\crefname{originallemma}{original lemma}{original lemmas}
\Crefname{originallemma}{Original lemma}{Original lemmas}
\makeatletter
\newlabel{morphisms-lemma-composition-finite-presentation}{{01TR}{source}{}{sourcepointer-2}{}}
\newlabel{morphisms-lemma-composition-finite-presentation@cref}{{[originallemma][0][]01TR}{[1][0][]source}{}{}{}}
\makeatother
\crefname{originallemma}{original lemma}{original lemmas}
\Crefname{originallemma}{Original lemma}{Original lemmas}
\makeatletter
\newlabel{morphisms-lemma-composition-flat}{{01U7}{source}{}{sourcepointer-3}{}}
\newlabel{morphisms-lemma-composition-flat@cref}{{[originallemma][0][]01U7}{[1][0][]source}{}{}{}}
\makeatother
\crefname{originallemma}{original lemma}{original lemmas}
\Crefname{originallemma}{Original lemma}{Original lemmas}
\makeatletter
\newlabel{morphisms-lemma-etale-flat}{{02GS}{source}{}{sourcepointer-4}{}}
\newlabel{morphisms-lemma-etale-flat@cref}{{[originallemma][0][]02GS}{[1][0][]source}{}{}{}}
\makeatother
\crefname{originallemma}{original lemma}{original lemmas}
\Crefname{originallemma}{Original lemma}{Original lemmas}
\makeatletter
\newlabel{morphisms-lemma-etale-locally-finite-presentation}{{02GR}{source}{}{sourcepointer-5}{}}
\newlabel{morphisms-lemma-etale-locally-finite-presentation@cref}{{[originallemma][0][]02GR}{[1][0][]source}{}{}{}}
\makeatother
\title{Postcomposition with an etale morphism preserves at-worst-nodal morphisms}
\author{The Stacks Project Authors}
\date{stacks-project:curves.tex:local-snapshot-20261012; GNU FDL 1.2 or later}
\begin{document}
\maketitle
\noindent\textit{Editorial note.} All delivered mathematical text is the author's own, copied from the pinned Stacks Project chapter file sources/stacks-curves/pinned-original-source.tex; no statement or proof step was rewritten, reordered or removed.\par\medskip
\noindent\textit{Editorial note.} The counted proof is the author's complete argument: the smooth-point case of the fibre statement is the standard fact that smoothness is stable under field extension (the source states it as 'the analogous result for smooth points' without a reference); the nodal case is Lemma lemma-node-over-separable-extension (in-source, navigation stub); the flatness and finite presentation steps cite the standard Morphisms lemmas (external tooling).\par\medskip
\noindent\textit{Editorial note.} The label lemma-nodal-family-postcompose-etale names the counted statement, and the proof environment immediately following it in the source is delivered verbatim as its proof.\par\medskip
\noindent\textit{Editorial note.} The counted statement is carried by the author's own lemma environment in the chapter source.\par\medskip
\noindent\textit{Allowed standard prerequisites (editorial).}\begin{itemize}
\item Fibral definition of at-worst-nodal morphisms. Same unit, the definition of at-worst-nodal morphisms of relative dimension 1. The proof reduces to the fibres because the at-worst-nodal condition on a flat locally finitely presented morphism is a condition on its fibres, by definition.
\item Nodal and smooth points under separable extension. Same unit, Lemma lemma-node-over-separable-extension. The nodal case of the fibre statement follows from Lemma lemma-node-over-separable-extension; the smooth-point case is the standard stability of smoothness under field extension.
\end{itemize}\medskip
\ifdefined\hypertarget\hypertarget{sourcepointer-1}{}\fi
\noindent\textit{Source reference.} Original reference 0CD6: sources/stacks-curves/pinned-original-source.tex, source lines 4684--4684. Prerequisite pointer invoked at proof line(s) 5035; it names a standard fact of the same chapter that the counted argument uses, and every step of the counted proof built on it is retained verbatim in the delivered body.\par\medskip
\ifdefined\hypertarget\hypertarget{sourcepointer-2}{}\fi
\noindent\textit{Source reference.} Original lemma 01TR: sources/stacks-chow/tags.txt, source lines 2384--2384. Prerequisite pointer invoked at proof line(s) 5031; it names a standard fact of a sibling chapter that the counted argument uses, and every step of the counted proof built on it is retained verbatim in the delivered body.\par\medskip
\ifdefined\hypertarget\hypertarget{sourcepointer-3}{}\fi
\noindent\textit{Source reference.} Original lemma 01U7: sources/stacks-chow/tags.txt, source lines 2400--2400. Prerequisite pointer invoked at proof line(s) 5032; it names a standard fact of a sibling chapter that the counted argument uses, and every step of the counted proof built on it is retained verbatim in the delivered body.\par\medskip
\ifdefined\hypertarget\hypertarget{sourcepointer-4}{}\fi
\noindent\textit{Source reference.} Original lemma 02GS: sources/stacks-chow/tags.txt, source lines 3239--3239. Prerequisite pointer invoked at proof line(s) 5030; it names a standard fact of a sibling chapter that the counted argument uses, and every step of the counted proof built on it is retained verbatim in the delivered body.\par\medskip
\ifdefined\hypertarget\hypertarget{sourcepointer-5}{}\fi
\noindent\textit{Source reference.} Original lemma 02GR: sources/stacks-chow/tags.txt, source lines 3238--3238. Prerequisite pointer invoked at proof line(s) 5029; it names a standard fact of a sibling chapter that the counted argument uses, and every step of the counted proof built on it is retained verbatim in the delivered body.\par\medskip
\begin{lemma}
\label{lemma-nodal-family-postcompose-etale}
Let $S' \to S$ be an \'etale morphism of schemes.
Let $f : X \to S'$ be at-worst-nodal of relative dimension $1$.
Then the composition $g : X \to S$
is at-worst-nodal of relative dimension $1$.
\end{lemma}

\begin{proof}
Observe that $g$ is flat and locally of finite presentation as
a composition of morphisms which are flat and locally of finite
presentation (use
Morphisms, Lemmas \ref{morphisms-lemma-etale-locally-finite-presentation},
\ref{morphisms-lemma-etale-flat},
\ref{morphisms-lemma-composition-finite-presentation}, and
\ref{morphisms-lemma-composition-flat}).
Thus it suffices to prove the fibres of $g$
have at-worst-nodal singularities.
This follows from Lemma \ref{lemma-node-over-separable-extension}
and the analogous result for smooth points.
\end{proof}

\end{document}
